% !TeX spellcheck = it_IT

\documentclass[format=169]{beamer}

\usepackage[utf8]{inputenc}
\usepackage[italian]{babel}
\usepackage{multicol}
\usetheme{metropolis}

\usepackage{listings}


\title{Esercitazioni di Informatica B}
\subtitle{Introduzione al C e costrutto while}

\author{Stefano Cereda\\
		stefano.cereda@polimi.it
	}
\date{8/10/2017}
\institute[PoliMi]{\vspace{0.5cm}\centering Politecnico di Milano \\ \vspace{0.2cm}
	\includegraphics[width=0.2\linewidth]{../logopolimi}}
%\logo{\includegraphics[width=15mm]{./logopolimi}}

\setbeamercovered{invisible}

\makeindex

\begin{document}
\begin{frame}
	\maketitle
\end{frame}

\begin{frame}{Indice}
\tableofcontents
\end{frame}

\section{Introduzione al C}

\begin{frame}{Esercizi di introduzione al C}
Si risolvano in C gli esercizi 1 e 6 di 'Introduzione agli algoritmi'.
\end{frame}

\begin{frame}
\frametitle{Esercizio 1 - Calcolatrice}
Scrivere un programma che, dopo aver ricevuto dall'utente due numeri interi ed un operatore matematico, stampi il risultato dell'operazione. Le operazioni da supportare sono: addizione, sottrazione, divisione, moltiplicazione e resto della divisione. Se l'utente inserisce un'operazione non supportata stampare un errore.
\end{frame}

\begin{frame}
\frametitle{Esercizio 2 - Statistiche}
Scrivere un programma che legga da tastiera 5 numeri e stampi a video il maggiore, la media e la radice quadrata della somma.

Per calcolare la radice quadrata si utilizzi la funzione \emph{sqrt} definita in \emph{math.h} (se utilizzate gcc aggiungete l'opzione \texttt{-lm})

Successivamente, si ripeta l'esercizio utilizzando il costrutto while.
\end{frame}

\begin{frame}
\frametitle{Esercizio 3 - Equazioni di secondo grado}
Scrivere un programma che, ricevuti tre numeri reali \emph{a,b,c}, calcoli e visualizzi le eventuali radici dell'equazione di secondo grado $a x^2 + b x +c = 0$
\end{frame}

\begin{frame}[fragile]
\frametitle{Esercizio 4 - Espressioni logiche}
\tiny
\alert{Tratto dal tema d'esame del 18/11/2008}
Si dica se le espressioni di seguito sono vere o false. Si dica inoltre se le espressioni sono sempre vere o sempre false nel caso in cui le variabili fossero inizializzate diversamente.

\begin{lstlisting}[basicstyle=\tiny]
void main()
{
int x=8,y=11;
char z='t';
// ESPRESSIONE 1
if ((y<7 && x>8) && (z>'a' || z<'z' ))
	printf("Espressione 1 vera\n");
else
	printf("Espressione 1 falsa\n");
// ESPRESSIONE 2
if ((-y==y) && (y >= 0))
	printf("Espressione 2 vera\n");
else
	printf("Espressione 2 falsa\n");
// ESPRESSIONE 3
if ((y>=10 || x<7) && (z=='c'))
{
	printf("Espressione 3 vera\n");
	// ESPRESSIONE 4
	if ((y>=10 || x<7) && (z<'c'))
		printf("Espressione 4 vera\n");
	else
		printf("Espressione 4 falsa\n");
	}
	else
		printf("Espressione 3 falsa\n");
}
\end{lstlisting}
\end{frame}

\begin{frame}
\frametitle{Esercizio 4 - Soluzione}
\begin{enumerate}
\item L'espressione è falsa per i valori dati a causa del valore di \emph{x}. Sarebbe però vera per \emph{x=9}, quindi non è sempre falsa.

\item L'espressione è falsa per i valori dati, ma sarebbe vera per \emph{y=0}.

\item L'espressione è falsa per i valori dati, ma sarebbe vera per \emph{x=6, z='c'}.

\item L'espressione 4, considerata l'espressione 3, sarà sempre falsa a causa di \emph{z}.
\end{enumerate}
\end{frame}


\iffalse
\begin{frame}
\frametitle{Esercizio 6 - Simulazione codice}
Dire cosa stampa il seguente codice:

TODO: LE ENUM NON LE HANNO ANCORA FATTE
%#include<stdio.h>
%const enum Alpha{
%	X,
%	Y=5,
%	Z
%}p=10;
%int main(){
%	enum Alpha a,b;
%	a= X;
%	b= Z;
%	printf("%d",a+b-p);
%	return 0;
%}
\end{frame}
\fi


\section{Costrutto while}
\begin{frame}
\frametitle{Esercizio 5 - Media di n numeri}
Scrivere un programma che legga una sequenza di numeri di lunghezza ignota. Ad ogni numero inserito il programma calcola e stampa la media dell'intera sequenza. La sequenza termina quando viene inserito un numero negativo.
\end{frame}

\begin{frame}
\frametitle{Esercizio 6 - Sottrazione ripetuta}
Si scriva un programma che legga da tastiera un numero intero \emph{n}, lo riduca ad un valore compreso fra 0 e 127 \emph{mediante sottrazione ripetuta} del valore 128, lo interpreti come carattere ASCII e lo stampi a video.
\end{frame}

\begin{frame}{Altri esercizi sul while}
Risolvere in C i rimanenti esercizi di 'Introduzione agli algoritmi'.
\end{frame}

\end{document}
